% !TEX program = pdflatex
\documentclass[11pt,a4paper]{article}

% =========================================================================
% GUÍA DE TRABAJO DEL INTEGRANTE
% Proyecto: CasaScore (AlquilaPuno)
% Gestión de Proyectos TICS (EST315) -- Semestre 2026-II -- Team Sobrevivientes
% =========================================================================

% --- PAQUETES DE CONFIGURACIÓN ---
\usepackage[utf8]{inputenc}
\usepackage[T1]{fontenc}
\usepackage[spanish,es-tabla]{babel}
\usepackage{geometry}
\geometry{
    top=3.6cm,
    bottom=2.3cm,
    left=2.2cm,
    right=2.2cm,
    % PARTE MODIFICABLE: headheight debe ser mayor que la altura del logo más alto (2.56cm = 73pt).
    headheight=80pt,
    headsep=15pt,
    footskip=20pt
}

\usepackage{graphicx}
\usepackage{fancyhdr}
\usepackage[table]{xcolor}
\usepackage{titlesec}
\usepackage{booktabs}
\usepackage{tabularx}
\usepackage{array}
\usepackage{amsmath}
\usepackage{enumitem}
\usepackage{microtype}
\usepackage{amssymb}
% --- RECURSOS VISUALES: GRÁFICOS Y CÓDIGO DE COLORES ---
\usepackage{tikz}
\usetikzlibrary{positioning,calc}
\definecolor{GrisFila}{RGB}{246, 248, 251}
\definecolor{CuadCerca}{RGB}{150, 30, 30}
\definecolor{CuadSatisf}{RGB}{180, 130, 20}
\definecolor{CuadInfo}{RGB}{10, 45, 90}
\definecolor{CuadMonit}{RGB}{110, 115, 125}
\newcommand{\etiqueta}[2]{{\setlength{\fboxsep}{2pt}\colorbox{#1}{\textcolor{white}{\scriptsize\bfseries #2}}}}
\newcommand{\cerca}{\etiqueta{CuadCerca}{Gestionar de cerca}}
\newcommand{\satisf}{\etiqueta{CuadSatisf}{Mantener satisfecho}}
\newcommand{\informado}{\etiqueta{CuadInfo}{Mantener informado}}
\newcommand{\monit}{\etiqueta{CuadMonit}{Monitorear}}
\usepackage{needspace}
\usepackage{longtable}
\definecolor{Rama1}{RGB}{10, 45, 90}
\definecolor{Rama2}{RGB}{40, 95, 150}
\definecolor{Rama3}{RGB}{180, 130, 20}
\definecolor{Rama4}{RGB}{30, 110, 60}
\definecolor{Rama5}{RGB}{110, 60, 140}
\definecolor{RaciR}{RGB}{150, 30, 30}
\definecolor{RaciA}{RGB}{10, 45, 90}
\definecolor{RaciC}{RGB}{180, 130, 20}
\definecolor{RaciI}{RGB}{140, 145, 155}
\definecolor{NivAlto}{RGB}{30, 110, 60}
\definecolor{NivMedio}{RGB}{180, 130, 20}
\definecolor{NivBajo}{RGB}{150, 30, 30}
\usetikzlibrary{patterns}
\definecolor{SevBaja}{RGB}{30, 110, 60}
\definecolor{SevMedia}{RGB}{180, 130, 20}
\definecolor{SevAlta}{RGB}{150, 30, 30}
\usetikzlibrary{arrows.meta,shapes.geometric,fit,backgrounds}
\usepackage{listings}
\lstset{basicstyle=\ttfamily\scriptsize, frame=single, rulecolor=\color{AzulUNA!40}, backgroundcolor=\color{GrisFila}, columns=fullflexible, keepspaces=true, breaklines=true, xleftmargin=4pt, xrightmargin=4pt, literate={á}{{\'a}}1 {é}{{\'e}}1 {í}{{\'i}}1 {ó}{{\'o}}1 {ú}{{\'u}}1 {ñ}{{\~n}}1 {Á}{{\'A}}1}

\usepackage{hyperref}

% Logos: funciona a 2 o 3 niveles de profundidad dentro de docs/latex/
\graphicspath{{../assets/img/}{../../assets/img/}{../../../assets/img/}}

% --- PALETA INSTITUCIONAL FIEI / UNA PUNO ---
\definecolor{AzulUNA}{RGB}{10, 45, 90}
\definecolor{DoradoFIEI}{RGB}{180, 130, 20}
\definecolor{GrisOscuro}{RGB}{50, 50, 50}
\definecolor{AzulClaro}{RGB}{228, 236, 246}

% --- CONFIGURACIÓN DE LISTAS ---
\setlist{topsep=2pt, itemsep=2pt, parsep=1pt}

% --- COLUMNAS DE TABLA ---
\newcolumntype{L}{>{\raggedright\arraybackslash}X}
\newcolumntype{P}[1]{>{\raggedright\arraybackslash}p{#1}}
\newcolumntype{C}[1]{>{\centering\arraybackslash}p{#1}}
\newcolumntype{Q}[1]{>{\raggedleft\arraybackslash}p{#1}}
\newcolumntype{R}{>{\raggedleft\arraybackslash}X}
\renewcommand{\arraystretch}{1.25}

% --- FORMATO DE ENLACES ---
\hypersetup{
    colorlinks=true,
    linkcolor=AzulUNA,
    urlcolor=AzulUNA,
    citecolor=AzulUNA,
    pdftitle={Guía De Trabajo Del Integrante -- CasaScore (AlquilaPuno)},
    pdfauthor={Team Sobrevivientes}
}

% --- FORMATO DE TÍTULOS ---
\titleformat{\section}{\color{AzulUNA}\normalfont\large\bfseries}{\thesection.}{0.7em}{}[\color{AzulUNA}\titlerule]
\titleformat{\subsection}{\color{AzulUNA}\normalfont\normalsize\bfseries}{\thesubsection.}{0.5em}{}
\titlespacing*{\section}{0pt}{9pt}{4pt}
\titlespacing*{\subsection}{0pt}{6pt}{2pt}

% --- ENCABEZADO Y PIE DE PÁGINA (fancyhdr) ---
\pagestyle{fancy}
\fancyhf{}

% ===== TAMAÑO DE LOS LOGOS (PARTE MODIFICABLE) =====
% logo_fiei.png  : height=1.45cm (lado izquierdo)
% logo_equipo.png: height=2.56cm (lado derecho; 25 % más grande porque la imagen tiene más borde blanco)
% Para cambiar el tamaño, modifique solo el valor height=... de cada \includegraphics.
% Si el logo más alto supera 2.6cm, aumente también headheight y top en la configuración de página.
\fancyhead[L]{%
    \includegraphics[height=1.45cm,keepaspectratio]{logo_fiei.png}%
}
\fancyhead[C]{%
    \scriptsize \textbf{\textcolor{AzulUNA}{UNIVERSIDAD NACIONAL DEL ALTIPLANO DE PUNO}}\\[2pt]
    \tiny \textcolor{GrisOscuro}{FACULTAD DE INGENIERÍA ESTADÍSTICA E INFORMÁTICA}\\[2pt]
    \scriptsize \textbf{\textcolor{AzulUNA}{Gestión de Proyectos TICS (EST315) \textbar\ Team Sobrevivientes}}\\[2pt]
    \footnotesize \textbf{\textcolor{DoradoFIEI}{GUÍA DE TRABAJO DEL INTEGRANTE}}%
}
\fancyhead[R]{%
    \includegraphics[height=2.56cm,keepaspectratio]{logo_equipo.png}%
}

\renewcommand{\headrulewidth}{0pt}
\renewcommand{\headrule}{}

\fancyfoot[L]{\scriptsize \textcolor{GrisOscuro}{CasaScore (AlquilaPuno)}}
\fancyfoot[C]{\small \thepage}
\fancyfoot[R]{\scriptsize \textcolor{GrisOscuro}{Puno, Perú -- Semestre 2026-II}}
\renewcommand{\footrulewidth}{0.4pt}

\begin{document}
\shorthandoff{<>}
\lstset{basicstyle=\ttfamily\fontsize{6.8}{8.2}\selectfont, breaklines=false}
\begin{center}
    {\Large \textbf{\textcolor{AzulUNA}{GUÍA DE TRABAJO DEL INTEGRANTE}}}\\[3pt]
    {\normalsize \textbf{Eliseo Tarqui Ajahuana -- Especialista Técnico}}\\[4pt]
    \footnotesize \textbf{Fecha:} jueves 05/11/2026 (semana 11) \quad \textbar\quad \textbf{Versión:} 1.0 \quad \textbar\quad \textbf{Encarga:} Clyde N. Paricahua Pari
\end{center}

\vspace{0.1cm}

\section{Propósito}
El equipo trabaja desde ahora en el repositorio \texttt{ProyectoTics}. Esta guía explica, desde cero, cómo preparar tu computadora, clonar el repositorio y subir tus documentos desde la terminal, cada uno con su rama, su commit y su Pull Request. El seguimiento de tus tareas se hace en el navegador, en el tablero \textit{CasaScore -- Tablero híbrido}.

\vspace{0.1cm}
\noindent\begin{tabularx}{\textwidth}{@{}P{3.6cm} L@{}}
\toprule
\textbf{Repositorio:} & \nolinkurl{https://github.com/xClaIDx/ProyectoTics} \\
\textbf{Tablero:} & \nolinkurl{https://github.com/users/xClaIDx/projects/1} \\
\textbf{Tu usuario:} & \texttt{Eliseo-Tarqui} \\
\textbf{Revisor:} & Clyde (\texttt{xClaIDx}) revisa, aprueba y fusiona cada Pull Request \\
\textbf{Plazo:} & PR 1 a 7 hasta el viernes 06/11; PR 8 el viernes 27/11 y PR 9 el viernes 04/12, en la fecha de cada documento \\
\bottomrule
\end{tabularx}

\section{Lo que Recibes}
\begin{enumerate}[leftmargin=*]
    \item \textbf{La carpeta \texttt{paquete\_eliseo}} (comprimida en un \texttt{.zip}), que Clyde te envía por WhatsApp. Contiene esta guía, la carpeta \nolinkurl{paquete_eliseo/docs/latex/} con tus 8 documentos (\texttt{.tex} y \texttt{.pdf}) y la carpeta \nolinkurl{paquete_eliseo/src/} con el código de la app móvil y de la app web, en las mismas rutas que tendrán en el repositorio.
    \item \textbf{La invitación al repositorio} \texttt{ProyectoTics}, que llega a tu correo de GitHub.
    \item \textbf{Tus tarjetas en el tablero}, ya asignadas a tu usuario.
\end{enumerate}

\Needspace{20\baselineskip}
\section{Preparación desde Cero (una sola vez)}

\subsection{Instalar Git y la terminal de GitHub}
{\small
\rowcolors{2}{white}{GrisFila}
\noindent\begin{tabularx}{\textwidth}{@{}P{2.2cm} L L@{}}
\toprule
\rowcolor{AzulClaro}
& \textbf{Windows} & \textbf{Mac} \\
\midrule
Git & Descarga \textbf{Git for Windows} desde \nolinkurl{https://git-scm.com/download/win} e instálalo con las opciones por defecto. & Abre la app \textbf{Terminal} y escribe \texttt{xcode-select --install}; acepta la instalación. \\
GitHub CLI & Descarga el instalador desde \nolinkurl{https://cli.github.com} e instálalo. & Descarga el instalador \texttt{.pkg} desde \nolinkurl{https://cli.github.com} e instálalo. \\
Terminal que usarás & \textbf{Git Bash} (búscalo en el menú Inicio). Todos los comandos de esta guía funcionan en Git Bash. & La app \textbf{Terminal}. \\
\bottomrule
\end{tabularx}
}

\vspace{0.1cm}
\noindent Cierra y vuelve a abrir la terminal, y comprueba que ambos programas responden con su número de versión:

\begin{lstlisting}
git --version
gh --version
\end{lstlisting}

\subsection{Configurar tu nombre}
Así cada commit queda a tu nombre. Usa el correo de tu cuenta de GitHub.

\begin{lstlisting}
git config --global user.name "Eliseo Tarqui Ajahuana"
git config --global user.email "tu-correo-de-github@ejemplo.com"
git config --global init.defaultBranch main
\end{lstlisting}

\Needspace{14\baselineskip}
\subsection{Iniciar sesión en GitHub desde la terminal}
\begin{lstlisting}
gh auth login
\end{lstlisting}

\noindent Responde así a cada pregunta (usa las flechas y Enter): \textbf{GitHub.com} $\rightarrow$ \textbf{HTTPS} $\rightarrow$ \textbf{Yes} (autenticar Git con tus credenciales) $\rightarrow$ \textbf{Login with a web browser}. La terminal muestra un código de 8 caracteres; pulsa Enter, se abre el navegador, pega el código y autoriza. Comprueba con:

\begin{lstlisting}
gh auth status
\end{lstlisting}

\Needspace{12\baselineskip}
\subsection{Aceptar la invitación al repositorio}
Desde el navegador, en el correo de GitHub (\textbf{View invitation} $\rightarrow$ \textbf{Accept invitation}) o en \nolinkurl{https://github.com/xClaIDx/ProyectoTics/invitations}. Sin este paso, la terminal rechazará tus envíos con un error de permiso.

\Needspace{16\baselineskip}
\subsection{Clonar el repositorio y preparar tu paquete}
\begin{enumerate}[leftmargin=*]
    \item Descomprime \texttt{paquete\_eliseo.zip} con doble clic y deja la carpeta \texttt{paquete\_eliseo} en tu \textbf{Escritorio}.
    \item En la terminal, ve al Escritorio y clona el repositorio:
\end{enumerate}

\begin{lstlisting}
cd ~/Desktop
git clone https://github.com/xClaIDx/ProyectoTics.git
cd ProyectoTics
ls docs/latex
\end{lstlisting}

\noindent Debes ver las carpetas \texttt{00\_comun}, \texttt{casascore} y \texttt{horus\_sentinel}. En Windows, si tu Escritorio está en OneDrive, usa \nolinkurl{cd ~/OneDrive/Escritorio} en lugar de \nolinkurl{cd ~/Desktop}. Desde ahora, todos los comandos se escriben \textbf{dentro de la carpeta} \texttt{ProyectoTics}; el paquete queda al lado, en \nolinkurl{../paquete_eliseo}.


\Needspace{12\baselineskip}
\section{Cómo se Simula el Historial}
Cada documento conserva la fecha que tiene escrita; no se modifica ningún archivo. GitHub registrará el día real en que lo subes. Para que se sepa cuándo debió entregarse, el mensaje del commit y el título del Pull Request terminan con la fecha prevista, por ejemplo:

\begin{lstlisting}
docs(casascore-requerimientos): requerimientos de ambos proyectos [fecha prevista 04/09/2026]
\end{lstlisting}

\Needspace{20\baselineskip}
\section{Tus Nueve Pull Requests}
Súbelos en este orden. Los documentos se copian desde \nolinkurl{../paquete_eliseo/docs/latex/} a \nolinkurl{docs/latex/}, y la app desde \nolinkurl{../paquete_eliseo/src/} a \nolinkurl{src/}.

\vspace{0.1cm}
{\small
\rowcolors{2}{white}{GrisFila}
\noindent\begin{tabularx}{\textwidth}{@{}C{0.6cm} P{4.6cm} L C{1.8cm}@{}}
\toprule
\rowcolor{AzulClaro}
\textbf{PR} & \textbf{Rama} & \textbf{Carpeta que subes} & \textbf{Fecha prevista} \\
\midrule
1 & \nolinkurl{docs/casascore-requerimientos} & \nolinkurl{casascore/02_requerimientos} y \nolinkurl{horus_sentinel/02_requerimientos} & 04/09/2026 \\
2 & \nolinkurl{docs/casascore-calidad} & \nolinkurl{casascore/11_plan_calidad} & 22/09/2026 \\
3 & \nolinkurl{docs/casascore-contratos} & \nolinkurl{casascore/13_plan_contratos} & 22/09/2026 \\
4 & \nolinkurl{docs/casascore-arquitectura} & \nolinkurl{casascore/14_arquitectura_tecnica} & 22/09/2026 \\
5 & \nolinkurl{feature/casascore-app-movil} & \nolinkurl{src/mobile} & 09/10/2026 \\
6 & \nolinkurl{docs/casascore-validacion} & \nolinkurl{casascore/17_informe_validacion} & 09/10/2026 \\
7 & \nolinkurl{feature/casascore-app-web} & \nolinkurl{src/web} & 06/11/2026 \\
8 & \nolinkurl{docs/casascore-auditoria} & \nolinkurl{casascore/24_informe_auditoria_calidad} & 27/11/2026 \\
9 & \nolinkurl{docs/casascore-manuales} & \nolinkurl{casascore/26_manuales} & 04/12/2026 \\
\bottomrule
\end{tabularx}
}

\vspace{0.1cm}
\noindent Los PR 8 y 9 corresponden a documentos con fecha 27/11 y 04/12/2026: súbelos en esas fechas, no antes.

\Needspace{18\baselineskip}
\section{Probar la App antes de Subirla}
Antes de los PR 5 y 7, comprueba que el código funciona en tu computadora. Necesitas \textbf{Node.js} en su versión LTS, que se descarga desde \nolinkurl{https://nodejs.org}. Las carpetas \texttt{node\_modules} que se crean al instalar no se suben: el archivo \texttt{.gitignore} del repositorio las excluye.

\begin{lstlisting}
node --version
cd ../paquete_eliseo/src/web
npm install
npm test
cd ../mobile
npm install
npx expo start
cd ../../../ProyectoTics
\end{lstlisting}

\noindent \texttt{npm test} debe terminar con todas las pruebas del motor aprobadas. \texttt{npx expo start} muestra un código QR para abrir la app móvil con Expo Go; ciérralo con Ctrl + C. Si algo falla, regístralo como tarjeta de tipo \textit{Defecto} antes de abrir el PR.


\Needspace{30\baselineskip}
\section{Ciclo Completo de un Pull Request}
Ejemplo con el PR 1. Cada bloque va precedido de lo que hace.

\vspace{0.1cm}
\noindent\textbf{1. Partir de la versión más reciente de \texttt{main}}
\begin{lstlisting}
git switch main
git pull
\end{lstlisting}

\noindent\textbf{2. Crear tu rama}
\begin{lstlisting}
git switch -c docs/casascore-requerimientos
\end{lstlisting}

\noindent\textbf{3. Copiar tus carpetas desde el paquete}
\begin{lstlisting}
cp -r ../paquete_eliseo/docs/latex/casascore/02_requerimientos docs/latex/casascore/
cp -r ../paquete_eliseo/docs/latex/horus_sentinel/02_requerimientos docs/latex/horus_sentinel/
\end{lstlisting}

\noindent\textbf{4. Revisar qué cambió}: deben aparecer en rojo las carpetas nuevas.
\begin{lstlisting}
git status
\end{lstlisting}

\noindent\textbf{5. Preparar y guardar el cambio (commit)}
\begin{lstlisting}
git add docs/latex/casascore/02_requerimientos
git add docs/latex/horus_sentinel/02_requerimientos
git commit -m "docs(casascore-requerimientos): requerimientos de ambos proyectos [fecha prevista 04/09/2026]"
\end{lstlisting}

\noindent\textbf{6. Enviar la rama a GitHub}
\begin{lstlisting}
git push -u origin docs/casascore-requerimientos
\end{lstlisting}

\noindent\textbf{7. Abrir el Pull Request con Clyde como revisor}
\begin{lstlisting}
gh pr create --base main --reviewer xClaIDx \
  --title "docs(casascore-requerimientos): requerimientos de ambos proyectos [fecha prevista 04/09/2026]" \
  --body "Fecha prevista 04/09/2026. Requerimientos de CasaScore y Horus Sentinel (.tex y .pdf). PR 1 de 9."
\end{lstlisting}

\noindent La terminal responde con el enlace del Pull Request. No lo fusiones tú: Clyde lo revisa y lo fusiona.

\vspace{0.1cm}
\noindent\textbf{8. Cuando Clyde lo fusione, volver a \texttt{main} y borrar la rama local}
\begin{lstlisting}
git switch main
git pull
git branch -d docs/casascore-requerimientos
\end{lstlisting}

\noindent No esperes a que fusione un PR para empezar el siguiente: el paso 1 siempre parte de \texttt{main}, y cada rama es independiente.

\Needspace{12\baselineskip}
\subsection{Si Clyde pide cambios}
Vuelve a tu rama, corrige el archivo, y envía un nuevo commit; el Pull Request se actualiza solo.
\begin{lstlisting}
git switch docs/casascore-requerimientos
git add docs/latex/casascore/02_requerimientos
git commit -m "fix(casascore-requerimientos): correcciones de la revision"
git push
\end{lstlisting}

\Needspace{12\baselineskip}
\subsection{Ver el estado de tus Pull Requests}
\begin{lstlisting}
gh pr status
gh pr list --author "@me"
gh pr view --web
\end{lstlisting}

\Needspace{16\baselineskip}
\section{Comandos de los PR 2 a 9}
Cada bloque se ejecuta igual que el ciclo anterior. Después de cada uno, cuando Clyde lo fusione, haz el paso 8 con el nombre de esa rama.

\Needspace{11\baselineskip}
\noindent\textbf{PR 2 -- Plan de calidad (22/09/2026)}
\begin{lstlisting}
git switch main && git pull
git switch -c docs/casascore-calidad
cp -r ../paquete_eliseo/docs/latex/casascore/11_plan_calidad docs/latex/casascore/
git add docs/latex/casascore/11_plan_calidad
git commit -m "docs(casascore-calidad): plan de calidad [fecha prevista 22/09/2026]"
git push -u origin docs/casascore-calidad
gh pr create --base main --reviewer xClaIDx --fill
\end{lstlisting}

\Needspace{11\baselineskip}
\noindent\textbf{PR 3 -- Plan de contratos y adquisiciones (22/09/2026)}
\begin{lstlisting}
git switch main && git pull
git switch -c docs/casascore-contratos
cp -r ../paquete_eliseo/docs/latex/casascore/13_plan_contratos docs/latex/casascore/
git add docs/latex/casascore/13_plan_contratos
git commit -m "docs(casascore-contratos): plan de contratos y adquisiciones [fecha prevista 22/09/2026]"
git push -u origin docs/casascore-contratos
gh pr create --base main --reviewer xClaIDx --fill
\end{lstlisting}

\Needspace{11\baselineskip}
\noindent\textbf{PR 4 -- Arquitectura técnica (22/09/2026)}
\begin{lstlisting}
git switch main && git pull
git switch -c docs/casascore-arquitectura
cp -r ../paquete_eliseo/docs/latex/casascore/14_arquitectura_tecnica docs/latex/casascore/
git add docs/latex/casascore/14_arquitectura_tecnica
git commit -m "docs(casascore-arquitectura): arquitectura tecnica [fecha prevista 22/09/2026]"
git push -u origin docs/casascore-arquitectura
gh pr create --base main --reviewer xClaIDx --fill
\end{lstlisting}

\Needspace{11\baselineskip}
\noindent\textbf{PR 5 -- App móvil del MVP de la Unidad I (09/10/2026)}
\begin{lstlisting}
git switch main && git pull
git switch -c feature/casascore-app-movil
mkdir -p src
cp -r ../paquete_eliseo/src/mobile src/
git add src/mobile
git commit -m "feat(casascore-app): app movil del MVP de la Unidad I [fecha prevista 09/10/2026]"
git push -u origin feature/casascore-app-movil
gh pr create --base main --reviewer xClaIDx --fill
\end{lstlisting}

\Needspace{11\baselineskip}
\noindent\textbf{PR 6 -- Informe de validación (09/10/2026)}
\begin{lstlisting}
git switch main && git pull
git switch -c docs/casascore-validacion
cp -r ../paquete_eliseo/docs/latex/casascore/17_informe_validacion docs/latex/casascore/
git add docs/latex/casascore/17_informe_validacion
git commit -m "docs(casascore-validacion): producto e informe de validacion [fecha prevista 09/10/2026]"
git push -u origin docs/casascore-validacion
gh pr create --base main --reviewer xClaIDx --fill
\end{lstlisting}

\Needspace{11\baselineskip}
\noindent\textbf{PR 7 -- App web (06/11/2026)}
\begin{lstlisting}
git switch main && git pull
git switch -c feature/casascore-app-web
mkdir -p src
cp -r ../paquete_eliseo/src/web src/
git add src/web
git commit -m "feat(casascore-web): app web con ranking, mapa y publicacion de cuartos [fecha prevista 06/11/2026]"
git push -u origin feature/casascore-app-web
gh pr create --base main --reviewer xClaIDx --fill
\end{lstlisting}

\Needspace{11\baselineskip}
\noindent\textbf{PR 8 -- Informe de auditoría y calidad (27/11/2026) -- súbelo el 27/11}
\begin{lstlisting}
git switch main && git pull
git switch -c docs/casascore-auditoria
cp -r ../paquete_eliseo/docs/latex/casascore/24_informe_auditoria_calidad docs/latex/casascore/
git add docs/latex/casascore/24_informe_auditoria_calidad
git commit -m "docs(casascore-auditoria): informe de auditoria y calidad [fecha prevista 27/11/2026]"
git push -u origin docs/casascore-auditoria
gh pr create --base main --reviewer xClaIDx --fill
\end{lstlisting}

\Needspace{11\baselineskip}
\noindent\textbf{PR 9 -- Manuales de usuario y técnico (04/12/2026) -- súbelo el 04/12}
\begin{lstlisting}
git switch main && git pull
git switch -c docs/casascore-manuales
cp -r ../paquete_eliseo/docs/latex/casascore/26_manuales docs/latex/casascore/
git add docs/latex/casascore/26_manuales
git commit -m "docs(casascore-manuales): manual de usuario y manual tecnico [fecha prevista 04/12/2026]"
git push -u origin docs/casascore-manuales
gh pr create --base main --reviewer xClaIDx --fill
\end{lstlisting}

\noindent La opción \texttt{--fill} usa el mensaje del commit como título y descripción del Pull Request.

\Needspace{24\baselineskip}
\section{Manejo del Proyecto en el Navegador}
El tablero se usa desde \nolinkurl{https://github.com/users/xClaIDx/projects/1}. Escribe \nolinkurl{assignee:Eliseo-Tarqui} en la barra de filtro para ver solo tus tarjetas.

\subsection{Qué haces con cada tarjeta}
{\small
\rowcolors{2}{white}{GrisFila}
\noindent\begin{tabularx}{\textwidth}{@{}P{4.4cm} C{1.9cm} L@{}}
\toprule
\rowcolor{AzulClaro}
\textbf{Tarjeta} & \textbf{Estado} & \textbf{Qué haces} \\
\midrule
Subir mis documentos y la app a ProyectoTics (9 PR) & Por hacer & Arrástrala a \textit{En Progreso} al empezar. Ábrela y marca la casilla de cada PR cuando lo crees. Pásala a \textit{Hecho} cuando Clyde fusione el PR 9. \\
Registro e inicio de sesión (Supabase) & En Progreso & Terminar el registro y el inicio de sesión con Ruth como apoyo en la seguridad por fila; fecha límite 18/11. \\
Punto de referencia en el mapa & Por hacer & Pásala a \textit{En Progreso} al empezar; fecha límite 13/11. \\
Perfil y mis cuartos; Pruebas de aceptación & Backlog & Se planifican en los sprints 9 y 11. \\
Informe de auditoría y calidad; Manuales de usuario y técnico & Backlog & Corresponden a los PR 8 y 9; actualízalos con los resultados reales antes de subirlos. \\
Motor CASA SCORE; Publicación de cuartos; Ranking y detalle; Modelo de datos y arquitectura; Diseño de pantallas & Hecho & Cuando se fusione el PR correspondiente, escribe en la tarjeta un comentario con el enlace del Pull Request. \\
\bottomrule
\end{tabularx}
}


\Needspace{12\baselineskip}
\subsection{Acciones en la tarjeta}
\begin{itemize}[leftmargin=*]
    \item \textbf{Cambiar de estado:} arrastra la tarjeta a otra columna, o ábrela y cambia el campo \textit{Status}.
    \item \textbf{Marcar un avance:} abre la tarjeta y marca la casilla de la lista en su descripción.
    \item \textbf{Registrar horas:} en el panel derecho de la tarjeta, suma el tiempo en \textit{Horas reales}, el mismo día y en decimales (media hora $=$ 0.5).
    \item \textbf{Dejar evidencia:} escribe un comentario abajo de la tarjeta con el enlace del Pull Request.
    \item \textbf{Avisar un bloqueo:} escribe el motivo como comentario y avisa a Clyde por WhatsApp el mismo día.
\end{itemize}

\Needspace{16\baselineskip}
\section{Comandos de Consulta Rápida}
{\small
\rowcolors{2}{white}{GrisFila}
\noindent\begin{tabularx}{\textwidth}{@{}P{5.6cm} L@{}}
\toprule
\rowcolor{AzulClaro}
\textbf{Comando} & \textbf{Para qué sirve} \\
\midrule
\texttt{pwd} & Ver en qué carpeta estás; debe terminar en \texttt{ProyectoTics}. \\
\texttt{git status} & Ver en qué rama estás y qué archivos cambiaron. \\
\texttt{git branch} & Listar tus ramas; la actual tiene un asterisco. \\
\texttt{git log --oneline -5} & Ver los últimos 5 commits. \\
\texttt{git switch main \&\& git pull} & Volver a \texttt{main} y traer lo último de GitHub. \\
\texttt{gh pr status} & Ver el estado de tus Pull Requests. \\
\bottomrule
\end{tabularx}
}

\Needspace{18\baselineskip}
\section{Si Algo Sale Mal}
{\small
\rowcolors{2}{white}{GrisFila}
\noindent\begin{tabularx}{\textwidth}{@{}P{5.2cm} L@{}}
\toprule
\rowcolor{AzulClaro}
\textbf{Mensaje o situación} & \textbf{Solución} \\
\midrule
\texttt{fatal: not a git repository} & No estás dentro de la carpeta; escribe \nolinkurl{cd ~/Desktop/ProyectoTics}. \\
\texttt{Permission denied} o error 403 al hacer \texttt{push} & Acepta la invitación al repositorio y vuelve a ejecutar \texttt{gh auth login}. \\
\texttt{Author identity unknown} & Falta configurar tu nombre y correo (sección 3.2). \\
\texttt{cp: cannot stat} & La carpeta \texttt{paquete\_eliseo} no está en el Escritorio, al lado de \texttt{ProyectoTics}. \\
\texttt{pathspec did not match} & La ruta está mal escrita; usa la tecla Tab para completarla. \\
\texttt{rejected (fetch first)} & Ejecuta \texttt{git pull} en tu rama y vuelve a hacer \texttt{git push}. \\
Hiciste un commit estando en \texttt{main} & No hagas \texttt{push}; avisa a Clyde y él te indica cómo moverlo a una rama. \\
\bottomrule
\end{tabularx}
}

\end{document}
